% !TEX root = ../main.tex
\chapter{A first session}\label{ch:first-session}
This chapter takes one recording from start to finish in about fifteen minutes. You find a
stretch of one channel in Explore, run a two-block chain over it in Analyse, and judge what
the chain found in Review. Each step names the control you press. Where the screenshot shows
a different recording or chain from the one in the steps, the caption says so; the controls
are the same.

\begin{plainterms}[The plan in one line]
Find a piece of the recording worth a look, ask the machine to mark every spike in it, then
go through its marks one by one and say which you agree with. It is the same loop as
proof-reading a spell-checker's suggestions: the program underlines, you accept or reject.
\end{plainterms}

\section{Before you start}\label{sec:first-session-before}
\begin{itemize}
  \item The app must be running and open in your browser. Someone who runs the project
    starts it with \code{webui\start.ps1} from the repository; it serves the
    app at \code{http://127.0.0.1:8765} (\code{CLAUDE.md}, ``Web UI'').
  \item The left edge of every page is the navigation rail: \ui{Explore}, \ui{Analyse},
    \ui{Discovery}, \ui{Models}, \ui{Review} and \ui{Library} at the top, \ui{Jobs} and
    \ui{Settings} at the foot. Across the top of the page a small path strip shows
    \ui{Explore}, \ui{Analyse}, \ui{Review} and lights the one you are on. Either takes you
    to a workspace.
  \item One recording, \code{M4_aug_concat_fs1.mat}, is \emph{held out}: it is kept unseen
    for the final test of the project, and every page refuses to open it. The header's
    \ui{M4 held out} chip is a reminder. Pick any other recording.
  \item A page marked with a dashed \ui{demo data} chip in its header is showing example
    data on part of the page, not your recording.
\end{itemize}

\section{Find a span in Explore}\label{sec:first-session-explore}
\begin{enumerate}
  \item Press \ui{Explore} in the rail. You land on \apppage{explore.corpus}: a coverage
    map with one row per channel and one column per time bin, shaded by how many annotations
    and detections fall in each bin.
  \item If the recording named at the top left is not the one you want, click its name to
    open the recordings menu and pick another, or step through recordings with the arrows
    beside it.
  \item Click a channel's row on the map to select it. The bar at the bottom names the
    channel and counts what is on it.
  \item Press \ui{Open \textit{channel}} at the bottom right (for example \ui{Open CH4\_A2}).
    Double-clicking a row does the same. You are now on \apppage{explore.signal}.
\end{enumerate}

\screenshot{pages/explore/explore.corpus--default}{The coverage map of a 16-channel,
721-hour recording. Darker bins hold more spans. CH4\_A2 is selected, and \ui{Open CH4\_A2}
at the bottom right opens it.}

The Signal page has three tiers on one time axis. The top tier, \textsf{CHANNEL}, shows the
whole channel with a blue box on it. The middle tier, \textsf{SPAN}, shows what is inside
that box at full detail. The bottom tier, \textsf{MOTIF}, shows one selected shape.

\begin{enumerate}[resume]
  \item Choose your span. Drag the blue box along the top tier, or drag its edges to widen
    or narrow it. In the middle tier, drag to pan and use the mouse wheel to zoom; the
    \ui{+} and \ui{−} buttons and the keys \texttt{+} and \texttt{−} zoom too, and
    \texttt{F} fits the whole channel. A span of an hour or two is a good size for a first
    chain.
  \item Read the line \ui{Selected span} in the action row under the tiers. It gives the
    start, the end and the length of what the middle tier shows. \emph{That} is the span
    you are about to send.
  \item Press \ui{Send span to Analyse} (or the key \texttt{a}). A message confirms the span
    was sent as the chain's source, and the app opens Analyse.
\end{enumerate}

\screenshot{pages/explore/explore.signal--default}{Explore's Signal page. The blue box on the
top tier sets the two-hour span drawn in the middle tier; \ui{Send span to Analyse} in the
bottom row sends exactly that span. This channel carries example data, hence the
\ui{demo data} chip.}

\begin{watchout}[Save span is not a send]
\ui{Save span} keeps only the tags and note you typed. \ui{Take span for Review} writes the
span to the database as your own annotation and puts it in a Review queue. Neither sends the
span to Analyse; only \ui{Send span to Analyse} does.
\end{watchout}

\section{Run a chain over it in Analyse}\label{sec:first-session-analyse}
You arrive on \apppage{analyse.chain}. The toolbar's source chip names your recording,
channel and span; the \textsf{Source} row draws the raw span. If the page still holds a
chain from an earlier visit, its blocks are listed under the source; delete them with each
row's bin icon, or add your blocks around them.

\begin{plainterms}[What the chain will do]
The first block, \blockref{preprocessing.detrend}, takes away the slow drift of the
baseline, the way you would level a photo taken on a slope before measuring anything in it.
The second, \blockref{detection.spike_v1}, then marks every point that stands clearly above
or below its own neighbourhood: a spike.
\end{plainterms}

\begin{enumerate}
  \item On an empty chain, a \ui{Start a chain} card offers \ui{Import a template},
    \ui{Add the first block}, and an example chain. Press \ui{Add the first block}. (On a
    chain that already has blocks, press the \ui{+ insert} pill where you want the block.)
  \item The \ui{Insert a stage} window opens. Its top line says which type of thing the new
    block must accept here: at the start of a chain, a \term{signal}{Signal}. Type
    \texttt{baseline} in \ui{search blocks}, or open the \ui{preprocess} tab, and click
    \ui{Baseline removal}. Blocks that do not fit at this point are greyed, each with the
    reason.
  \item Press \ui{+ Insert}. The window closes and the chain gains row \textsf{01 Baseline
    removal}, marked \ui{new}. Its caption reads its main setting: a 600-second rolling mean
    is subtracted.
  \item Press \ui{+ insert · end of chain} below it, search for \texttt{spike}, click
    \ui{Spike detection (v1)} under the \ui{detect} tab, and press \ui{+ Insert}. Row
    \textsf{02} appears.
  \item Press \ui{Run chain} at the top right. Each row fills in as its step finishes:
    the baseline row draws the signal before (grey) and after (blue); the spike row draws its
    spans over the signal, one per spike, and counts them, for example ``\textit{n} spans ·
    mean \textit{t} s · \textit{u} up, \textit{v} down''. The footer gives the run its number and says how many detections it wrote.
\end{enumerate}

\screenshot{pages/analyse/analyse.chain--2-insert-stage}{The \ui{Insert a stage} window.
The three pills at the top say what the block before gives and what the block after needs;
\ui{Fits here} lists the blocks that match, and the rest are greyed with their reason.}

\screenshot{pages/zz_analyse_run/analyse.chain--fixup-h-thumbnails-mp_threshold}{A chain
after a run (here a baseline removal, a matrix profile and a threshold, not the two blocks
in the steps). Every row shares the time axis at the bottom; the footer reads
``8 written to detections'' and offers \ui{→ Pass 8 to Review}.}

\begin{watchout}[If the spike row finds nothing, or far too much]
Spike detection's threshold, \emph{delta\_mv}, is in millivolts; the block converts it by
the unit the recording declares. Its default, 1\,mV, is a placeholder chosen for one
recording's noise (\cref{block:detection.spike_v1}). \unverified{Whether the default finds
spikes on your span depends on the recording.} To change it, press the row's settings icon
(the first of the four small icons) to open the block's page, lower \emph{delta\_mv} under
\ui{Parameters} if it finds nothing or raise it if it fires on noise, and press \ui{Re-run
from 02}. The readout beside the settings turns amber when the threshold is close to the
span's own wobble. \ui{full chain} takes you back. If the chain row says ``read as volts'',
the recording's unit is not declared: declare it in \apppage{settings.datasets}.
\end{watchout}

The spike detector draws one short stretch around each spike, with the spike's own moment
marked inside it. Each stretch is one detection, so each is one item to judge in Review.

\section{Judge what it found in Review}\label{sec:first-session-review}
\begin{enumerate}
  \item In the chain's footer press \ui{→ Pass \textit{N} to Review}. The app makes a Review
    queue over the \textit{N} detections the run wrote, names it after the chain and the run
    number, and opens it on \apppage{review.inspector}. The button stays grey until the run
    has finished, and greys again if you change the chain without re-running it.
  \item The page shows one candidate at a time. \ui{Candidate in context} draws the detection
    as a blue band on the surrounding signal; the \ui{padding} buttons (\ui{±30 s},
    \ui{±120 s}, \ui{±300 s}) set how much surrounds it. The counter at the top says how many
    are judged out of how many.
  \item Give a verdict by pressing its key or its card under \ui{Verdict}:
    \texttt{I} \ui{interesting}, \texttt{N} \ui{not interesting}, \texttt{A} \ui{artifact}
    (a recording fault, kept out of training), \texttt{U} \ui{unsure}. \texttt{Space}
    \ui{skip} moves on without writing anything. With \ui{auto-advance} on, the next
    candidate opens by itself.
  \item Made a mistake? \texttt{Ctrl}+\texttt{Z} undoes the last verdict, in the database
    too. The arrow keys move to the next and previous candidate.
  \item Work through a dozen or so. The chip \ui{writes adjudications} in the toolbar is
    the app telling you where your verdicts go: to the human-verdict table, never into the
    machine's rows (\cref{sec:core-ideas-detections}).
\end{enumerate}

\screenshot{pages/review/review.inspector--1-candidate (deep link, live detection)}{One
candidate in Review (from a drop detector's queue, not the spike chain). The candidate is
the blue band; the verdict cards along the bottom each show their key.}

\begin{watchout}[S is not ``skip'']
\texttt{S} gives the verdict \ui{seed}, which promotes the candidate to a new exemplar in
the Library. Use it only for a shape you want to search for again. To pass over a
candidate, press \texttt{Space}.
\end{watchout}

\section{Where to go next}\label{sec:first-session-next}
You have now done the whole loop once. From here:
\begin{itemize}
  \item Press \ui{Save template} in Analyse's toolbar to keep your chain. A
    \term{template}{template} is the chain without the recording and span, ready to apply
    elsewhere (\cref{ch:analyse}).
  \item Run the template across many channels, with a check against chance, in Discovery
    (\cref{ch:discovery}).
  \item Learn the words this chapter used without stopping to explain them:
    \cref{ch:core-ideas}.
  \item Read each page in full: Explore (\cref{ch:explore}), Analyse (\cref{ch:analyse}),
    Review (\cref{ch:review}).
\end{itemize}

\begin{underthehood}[What the clicks call]
\ui{Send span to Analyse}: \code{webui/client/src/explore/SignalPage.tsx}, \code{send},
which sets the chain's source from the viewport. \ui{Run chain}:
\code{webui/client/src/analyse/ChainPage.tsx}, \code{run}; the run itself is
\code{Working/execution.py}, \code{execute_recipe}, which writes a terminal span set as
detections. \ui{Pass N to Review}: \code{webui/client/src/analyse/sendToReview.ts},
\code{useSendRunToReview} (\code{POST /api/review/queues}). Verdict keys:
\code{webui/client/src/review/keys.ts}.
\end{underthehood}
